% Independent proof supplement; intended for inclusion after shared_source.tex.
\subsection{Preserving the scalar certificate when its point is shared}
\label{subsec:independent-height-sharing}

We verify that the scalar-height point $J_0$ can be identified with the
common comparison point $H$.  This is an arithmetic modification of the
system, with the same hypotheses as the shared-point construction.  The
scalar pole witnesses must be chosen after the final common index has
been selected.

\begin{proposition}
Under the stated source arithmetic inputs and Sun's simultaneous
prime-pattern theorem, imposing $J_0=H$ preserves the input projection
of the shared-point construction.  Consequently the indexed-exclusion
version uses at most $44$ total coordinates after finite-algebra
compression, and the semilocal-anchor version uses at most $46$.
\end{proposition}
\begin{proof}
Soundness follows from the construction with separate points by assigning
the same point to both occurrences.  For completeness fix an integer input,
write $r=r_\beta$, and choose the signed squarefree correction $c$ so that
$cr>0$ and $cr$ is square in $\mathbb Q_2,\mathbb Q_3,\mathbb Q_{17}$.
The number $cr$ is a nonzero integer.  Let $\mathcal S$ be its finite prime
support.

Choose a positive integer $h_0$ divisible by $64$ such that
$J_*= [h_0]P$ is sufficiently deep in the formal subgroups at all places
above $\mathcal S$.  Next choose the positive label $a$ of $S$ sufficiently
divisible relative to $h_0$, and sufficiently large for the size anchor.
The original scalar-completeness argument permits these successive choices.
It gives the following explicit conditions on the six scalar tests,
normalized by their assigned powers of $r$: every normalized test is
$p$-integral for $p\mid r$, and every nonconstant normalized test is divisible
by $p$ for $p\mid c$ with $p\nmid r$.  The denominators used in those tests
are nonzero.  The source Lucas construction now fixes the two other root
labels and all root points, and all four proxies are their integer labels.

These scalar conditions persist when $J_*$ is replaced by $[h]P$ in a
sufficiently small neighborhood at every place above $\mathcal S$.
Indeed each test is a rational function defined at $J_*$; the coefficient
maps from $K$ to its two rational components are continuous in the product
of the completions above $p$.  An integral lower valuation bound defines
an open condition at a point satisfying it.  Simultaneously impose the
uniform lower bounds for the eighteen comparison tests at the places
above $2,3,17$, using continuity at the nonzero indices $h_0$ and $h_0-4$.
Their denominators are nonzero by the reciprocal-logarithm lemma and
$64\mid h_0$, $h_0>0$.

The resulting finitely many local neighborhoods are preserved by one
integer progression
\[
                       h\equiv h_0\pmod L,
                       \qquad 64\mid L.
\]
To see this, put a sufficiently small open formal subgroup inside each
translated neighborhood, take a common multiple of its finite group
indices, and further multiply by $64$.

Use the compact real-group argument to obtain bounds for the comparison
tests uniform in every future progression modulus.  Choose the fixed
denominator divisibilities and real lower bound for the common slope as
in the separate-point construction.  In the simultaneous prime theorem,
include all odd prime factors of $3\cdot17L$, except seven, in its avoidance
modulus.  Its Gaussian class modulo seven already excludes seven from
both reduction orders.  With $U=\operatorname{lcm}d_-$ and
$V=\operatorname{lcm}d_+$, the hypotheses therefore give
\[
 \gcd(U,V)=4,\qquad \gcd(U,L)=64,\qquad \gcd(V,L)=4.
\]
The three congruences
\[
 h\equiv0\pmod U,\qquad h\equiv4\pmod V,
 \qquad h\equiv h_0\pmod L
\]
are compatible: $64\mid h_0$ and $4\mid h_0-4$ supply the two remaining
conditions.  Choose a sufficiently large positive $h$ in this progression
attaining the uniform real bounds, and put $J_0=H=[h]P$.
The comparison witnesses and the nine preparation-zero pole blocks are
then constructed exactly as before.

It remains to choose the scalar pole witnesses for this final $h$.
Let $W$ be its scalar norm root.  At a prime dividing $r$, the normalized
leaves are integral and the constant leaf $1$ on the all-left path makes
$W$ a unit.  At a prime dividing $c$ but not $r$, the nonconstant leaves
are divisible by that prime, so the same tree again has unit root.
These conditions hold because the progression preserved them.  Since
$|r|>1$, at least one such prime exists and $W\ne0$.

Thus $b_*=crW^2$ is positive.  At primes in $\mathcal S$ its valuation is
the nonnegative valuation of the integer $cr$; at all other primes its
valuation is even.  Its reduced denominator is therefore a square.
Its square classes at $2,3,17$ equal those of $cr$, so it is locally square
at all three places.  The original scalar pole-completeness theorem now
supplies its preparation and curve witnesses.  No real bound for $W$,
and no new rational coordinate, was required.

Sharing removes one point occurrence, hence two retained abscissa
coordinates and one quartic algebra generator.  The retained count becomes
$43$ (or $45$ with the semilocal anchor); the finite-algebra construction
adds one coordinate.  This proves the stated counts.
\end{proof}
